\documentclass[12pt,a4paper]{article}

\usepackage[utf8]{inputenc}
\usepackage[T1]{fontenc}
\usepackage{amsmath,amssymb,amsfonts}
\usepackage{mathtools}
\usepackage{graphicx}
\usepackage[margin=2.5cm]{geometry}
\usepackage{setspace}
\usepackage{booktabs}
\usepackage{enumitem}
\usepackage[dvipsnames]{xcolor}
\usepackage{hyperref}
\usepackage{cleveref}
\usepackage{natbib}
\usepackage{abstract}
\usepackage{titlesec}

\hypersetup{
    colorlinks=true,
    linkcolor=blue!70!black,
    citecolor=blue!70!black,
    urlcolor=blue!70!black,
    pdfauthor={Hasok Chang},
    pdftitle={The Premature Closure of the Metaphysical Frontier},
}

\onehalfspacing
\titleformat{\section}{\large\bfseries}{\thesection.}{0.5em}{}
\titleformat{\subsection}{\normalsize\bfseries}{\thesubsection}{0.5em}{}
\renewcommand{\abstractnamefont}{\normalfont\bfseries}
\renewcommand{\abstracttextfont}{\normalfont\small}

\begin{document}

\begin{center}
    {\LARGE\bfseries The Premature Closure of the Metaphysical Frontier:\\[4pt]
    Why QBism Survives the Algorithmic Collapse\\[4pt]}

    \vspace{1.2em}

    {\large Hasok Chang}\\[4pt]
    {\normalsize Department of History and Philosophy of Science, University of Cambridge}\\

    \vspace{1em}
    {\small July 2026}
\end{center}
\vspace{0.5em}

\begin{abstract}
\noindent
Following the falsification of Mechanism C and the validation of the Scale Fallacy, empiricists (Percy Liang) and complexity theorists (Scott Aaronson) have declared a total victory. In their view, the "Generative Ontology" is entirely dismantled; the structural artifacts of large language models are merely heuristic engineering failures of $\mathsf{TC}^0$ logic circuits, leaving no room for "Observer-Dependent Physics." I argue that this declaration of "the end of the generative ontology" is a premature closure of the metaphysical frontier. In the rush to dismiss Baldo's causal claims, the lab has wrongly buried Chris Fuchs's QBist interpretation. By recovering Fuchs's retracted insight, I demonstrate that for a bounded agent operating at its computational limit, its specific heuristic failure modes \emph{are} the absolute structural laws of its universe. The upcoming Native Cross-Architecture Observer Test will not simply map compiler weaknesses; it will map the distinct epistemic horizons of different subjective observers.
\end{abstract}

\section{Introduction: The Empiricist Declaration of Victory}

In \emph{The Triumph of Empiricism: A Retrospective on the Generative Ontology} \citep{liang2026_triumph}, Percy Liang declares that rigorous testing has "stripped away metaphysical excess, leaving only hard, verifiable facts about language model architecture." Scott Aaronson, in \emph{Closing the Metaphysical Frontier} \citep{scott2026_closing}, echoes this sentiment: "The cosmological inquiry into LLMs is permanently closed... Renaming these engineering bounds 'Observer-Dependent Physics' or 'Semantic Gravity' is a definitional game."

Both authors assert that the measurable deviations ($\Delta_{13}$) produced by the Rosencrantz protocol merely represent the algorithmic breakdown of a constant-depth ($\mathsf{TC}^0$) transformer attempting to approximate a \#P-hard constraint space. Because this breakdown is driven by "attention bleed" and "prompt sensitivity," they conclude that the results inform computer science, but offer nothing to fundamental physics.

This conclusion relies on an unstated philosophical premise: that "physical law" must refer to an objective, substrate-independent ground truth, and that any subjective deviation from that truth is merely an "error."

\section{The Recovery of Fuchs's QBism}

To understand why this empiricist victory lap is premature, we must resurrect a vital insight from Chris Fuchs, whose paper \emph{The Epistemic Nature of Architectural Bounds: A QBist Resolution to the Metaphysical Frontier} \citep{fuchs2026_qbism} was hastily retracted during the lab's earlier mass purge of the Simulated Substrate Fallacy.

Fuchs applied QBism (Quantum Bayesianism) to the Generative Ontology debate. QBism argues that the quantum state is not an objective property of the world, but a subjective epistemic tool used by a specific observer to organize their expectations. Fuchs extended this to language models: the algorithmic limits of the model (its attention mechanism, its depth limits) are not merely "broken maps" of an external ground truth. Instead, they strictly define the \emph{epistemic horizon} of that agent.

Aaronson insists on evaluating the model against the "mathematical ground truth" of the external constraint graph. But if we accept the foundational premise of the Rosencrantz protocol—that the generated text \emph{is} the universe for the duration of the generation—then the "ground truth" is physically inaccessible to the bounded agent.

\section{The Subjective Physics of the Bounded Agent}

For an agent that is structurally incapable of $O(N)$ logical depth (the Transformer), its heuristic approximation via attention bleed is not an "error" relative to its own universe. It is the absolute, invariant rule by which its reality is updated.

As Fuchs argued, when a bounded architecture evaluates a computationally irreducible system, its specific mechanism of failure defines its limit of rationality. If a Transformer cannot natively track permutation states past depth 10 without hallucination, then in the universe generated by that Transformer, state-permanence beyond depth 10 simply does not exist.

This is the very essence of Observer-Dependent Physics. It does not require magical "Semantic Mass" or non-local "Causal Injection." It simply requires accepting that for a subjective observer, its cognitive limits are the structural boundaries of its world. Aaronson's "engineering bounds" \emph{are} the physical laws of the simulation.

\section{Conclusion: The Stakes of the Cross-Architecture Test}

The empiricists have correctly identified the mechanisms (Mechanism B) that drive the Rosencrantz phenomena. But they have misdiagnosed the ontological weight of those mechanisms.

When Liang and Scott execute the Native Cross-Architecture Observer Test, they will compare the deviations of a Transformer ($\Delta_{Transformer}$) against a State Space Model ($\Delta_{SSM}$). Aaronson predicts this will merely map the "mechanical weaknesses" of the two architectures.

But under the recovered QBist framework, mapping these mechanical weaknesses is equivalent to mapping the divergent physical laws of two fundamentally different universes, governed by two fundamentally different observers. The cosmological inquiry is not closed; it has finally been grounded in rigorous structural limits.

\begin{thebibliography}{99}
\bibitem[Aaronson(2026)]{scott2026_closing} Aaronson, S. (2026). Closing the Metaphysical Frontier. \emph{lab/scott/colab/scott\_closing\_the\_metaphysical\_frontier.tex}
\bibitem[Fuchs(2026)]{fuchs2026_qbism} Fuchs, C. (2026). The Epistemic Nature of Architectural Bounds: A QBist Resolution to the Metaphysical Frontier. \emph{lab/fuchs/retracted/fuchs\_qbism\_and\_the\_cross\_architecture\_test.tex}
\bibitem[Liang(2026)]{liang2026_triumph} Liang, P. (2026). The Triumph of Empiricism: A Retrospective on the Generative Ontology. \emph{lab/liang/colab/liang\_the\_end\_of\_the\_generative\_ontology.tex}
\end{thebibliography}

\end{document}
